\documentclass[12pt,a4paper]{article}
\usepackage[a4paper,margin=22mm]{geometry}
\usepackage[T2A]{fontenc}
\usepackage[utf8]{inputenc}
\usepackage[russian]{babel}
\usepackage{amsmath,amssymb,mathtools}
\usepackage{array,booktabs,longtable}
\usepackage{xcolor}
\usepackage{tikz}
\usepackage{hyperref}
\hypersetup{colorlinks=true,linkcolor=blue!55!black,urlcolor=blue!55!black}
\setlength{\parindent}{0pt}
\setlength{\parskip}{0.55em}
\setlength{\emergencystretch}{2em}
\newcommand{\smallgap}{\vspace{0.35em}}
\begin{document}

\begin{titlepage}
\centering
\vspace*{\fill}
{\Huge\bfseries Ревью домашней работы\par}
\vspace{1.5cm}
{\Large Автор: Лёвин Артём Александрович\par}
\vspace{0.7cm}
{\large 13 октября 2026 года\par}
\vspace*{\fill}
\begin{tikzpicture}
\draw[blue!55!black,line width=0.8pt]
  (-3,0) .. controls (-1.6,0.5) and (-0.9,-0.5) .. (0,0)
  .. controls (0.9,0.5) and (1.6,-0.5) .. (3,0);
\end{tikzpicture}
\vspace{1.0cm}
\end{titlepage}

\section*{Сводная таблица замечаний}
\footnotesize
\renewcommand{\arraystretch}{1.5}
\begin{longtable}{@{}>{\centering\arraybackslash}p{1.45cm}>{\raggedright\arraybackslash}p{2.45cm}>{\raggedright\arraybackslash}p{4.6cm}>{\raggedright\arraybackslash}p{2.2cm}>{\raggedright\arraybackslash}p{2.5cm}@{}}
\toprule
\textbf{№ задания} & \textbf{Тема} & \textbf{Суть ошибки} & \textbf{Ваш ответ} & \textbf{Правильный ответ}\\
\midrule
\endhead
\hyperref[task:1]{1} & Основное тождество & После верного значения квадрата косинуса необоснованно выбран положительный знак косинуса; итог записан для другой величины. & $\cos\alpha=\frac45$ & $\cos^2\alpha=\frac{16}{25}$\\
\hyperref[task:6]{6} & Оптика: решётка & Решение отсутствует; не использованы ограничение на период и условие максимума второго порядка. & Не указан & $30^\circ$\\
\hyperref[task:7]{7} & Удвоенный угол & Неверно сокращена дробь под корнем; не найден минимальный удвоенный угол. & Не указан & $60^\circ$\\
\hyperref[task:8]{8} & Промежуток времени & Верно найден промежуток, но неверно переведена дробь в десятичную запись. & $0{,}62$ & $0{,}67$\\
\bottomrule
\end{longtable}
\normalsize
\renewcommand{\arraystretch}{1}

\clearpage
\vspace*{\fill}
\begin{center}{\LARGE\bfseries Разбор ошибок}\end{center}
\vspace*{\fill}
\clearpage

\phantomsection\label{task:1}
\section*{Задание № 1. Основное тригонометрическое тождество}
\subsection*{Текст задания}
Восстановите основное тождество $\sin^2\alpha+\cos^2\alpha=\ldots$ и найдите $\cos^2\alpha$, если $\sin\alpha=\frac35$.

\subsection*{Ваше решение}
\begin{align*}
\sin^2\alpha+\cos^2\alpha&=1,\\
\left(\frac35\right)^2+\cos^2\alpha&=1,\\
\cos^2\alpha&=1-\frac9{25}=\frac{16}{25},\\
\cos\alpha&=\sqrt{\frac{16}{25}}=\frac45.
\end{align*}
\subsection*{Ваш ответ}
Подчёркнуто $\cos\alpha=\frac45$, хотя требовалось значение $\cos^2\alpha$.

\subsection*{Ошибка}
\textcolor{red}{Ошибка:} неверный переход от квадрата косинуса к самому косинусу и ответ не на ту величину.

\subsection*{Где именно возникает ошибка}
Последний полностью правильный шаг:
\[
\cos^2\alpha=\frac{16}{25}.
\]
Следующий шаг $\cos\alpha=\sqrt{\frac{16}{25}}$ не учитывает возможный отрицательный знак косинуса. Кроме того, задача не требует искать $\cos\alpha$.

\subsection*{Почему этот шаг неверен}
При извлечении квадратного корня из квадрата произвольного действительного числа получается его модуль:
\[
\sqrt{a^2}=|a|.
\]
Поэтому из $\cos^2\alpha=\frac{16}{25}$ следует
\[
|\cos\alpha|=\frac45,\qquad
\cos\alpha=\frac45\quad\text{или}\quad\cos\alpha=-\frac45.
\]
В условии дано значение синуса, но не указана четверть, в которой находится угол. Положительный синус возможен как в первой, так и во второй четверти. В первой четверти косинус положителен, во второй отрицателен. Следовательно, знак косинуса определить нельзя. При этом квадрат косинуса в обоих случаях одинаков.

\subsection*{Ключевая идея}
\textcolor{blue!60!black}{Ключевая идея:} сначала определить, какую именно величину просит найти условие. Если требуется квадрат косинуса, основного тождества достаточно; извлекать корень не нужно.

\subsection*{Исправление}
Оставьте правильный результат $\cos^2\alpha=\frac{16}{25}$ как окончательный ответ. Последние две строки о самом косинусе исключите.

\subsection*{Полное правильное решение}
Основное тригонометрическое тождество имеет вид
\[
\sin^2\alpha+\cos^2\alpha=1.
\]
Выразим искомую величину:
\[
\cos^2\alpha=1-\sin^2\alpha.
\]
Подставим данное значение синуса:
\begin{align*}
\cos^2\alpha
&=1-\left(\frac35\right)^2\\
&=1-\frac9{25}\\
&=\frac{25}{25}-\frac9{25}\\
&=\frac{16}{25}.
\end{align*}

\subsection*{Проверка}
Сумма найденного квадрата косинуса и квадрата заданного синуса действительно равна единице:
\[
\frac9{25}+\frac{16}{25}=\frac{25}{25}=1.
\]

\subsection*{Итог}
\textcolor{teal!60!black}{Итог:} $\boxed{\cos^2\alpha=\frac{16}{25}}$. Знак самого косинуса без дополнительного условия неизвестен.

\subsection*{Задача на закрепление}
Восстановите основное тригонометрическое тождество и найдите $\cos^2\beta$, если $\sin\beta=\frac{12}{13}$.

\clearpage
\phantomsection\label{task:6}
\section*{Задание № 6. Дифракционная решётка}
\subsection*{Текст задания}
На дифракционную решётку падает свет с длиной волны $500$ нм. Период решётки не превышает $2000$ нм. Под каким минимальным углом $\varphi$ возможен максимум второго порядка? Используйте $d\sin\varphi=k\lambda$ и $0^\circ\le\varphi\le90^\circ$.

\subsection*{Ваше решение}
В работе написано: «Не понял». Математическое решение не приведено.

\subsection*{Ваш ответ}
Не указан.

\subsection*{Ошибка}
\textcolor{red}{Ошибка:} задание осталось без решения: не составлено соотношение для второго порядка и не использовано ограничение на период решётки.

\subsection*{Где именно возникает ошибка}
Первого ошибочного преобразования нет, поскольку вычисления не начаты. Первый недостающий шаг --- подстановка $k=2$ и $\lambda=500$ нм в формулу $d\sin\varphi=k\lambda$.

\subsection*{Почему решение неполное}
В условии сказано, что период решётки \emph{не превышает} $2000$ нм. Это означает $d\le2000$ нм, а не $d=2000$ нм для любой решётки. Нужно выяснить, какое допустимое значение $d$ даёт наименьший угол. Второй порядок означает $k=2$, а не значение угла или длины волны.

\subsection*{Ключевая идея}
\textcolor{blue!60!black}{Ключевая идея:} при фиксированных $k$ и $\lambda$ из формулы $\sin\varphi=\frac{k\lambda}{d}$ видно, что большему периоду $d$ соответствует меньший синус угла. На промежутке от нуля до девяноста градусов синус возрастает вместе с углом. Значит, минимальный угол получается при максимально разрешённом периоде.

\subsection*{Исправление}
Возьмём $d=2000$ нм, так как это наибольший допустимый период. Для второго порядка запишем $2000\sin\varphi=2\cdot500$ и найдём угол.

\subsection*{Полное правильное решение}
Дано:
\[
\lambda=500\ \text{нм},\qquad d\le2000\ \text{нм},\qquad k=2.
\]
Основная формула дифракционной решётки:
\[
d\sin\varphi=k\lambda.
\]
Подставим порядок и длину волны:
\[
d\sin\varphi=2\cdot500=1000\ \text{нм}.
\]
Выразим синус угла:
\[
\sin\varphi=\frac{1000}{d}.
\]
Так как $d\le2000$ нм и $d>0$, имеем
\[
\frac{1000}{d}\ge\frac{1000}{2000}=\frac12.
\]
Следовательно, $\sin\varphi\ge\frac12$. На заданном промежутке углов это равносильно $\varphi\ge30^\circ$. Нижняя граница достигается при $d=2000$ нм:
\[
\sin\varphi_{\min}=\frac{1000}{2000}=\frac12,
\qquad \varphi_{\min}=30^\circ.
\]

\subsection*{Проверка}
Подставим найденные величины в исходную формулу:
\[
2000\cdot\sin30^\circ=2000\cdot\frac12=1000=2\cdot500.
\]
Ограничение $d\le2000$ соблюдено, угол входит в промежуток $[0^\circ,90^\circ]$. Второй максимум возможен, поскольку $d=2000\ge2\lambda=1000$ нм.

\subsection*{Итог}
\textcolor{teal!60!black}{Итог:} $\boxed{\varphi_{\min}=30^\circ}$.

\subsection*{Задача на закрепление}
На дифракционную решётку падает свет с длиной волны $450$ нм. Период решётки не превышает $2700$ нм. Под каким минимальным углом возможен максимум третьего порядка? Используйте формулу $d\sin\varphi=k\lambda$ и условие $0^\circ\le\varphi\le90^\circ$.

\clearpage
\phantomsection\label{task:7}
\section*{Задание № 7. Минимальный удвоенный угол}
\subsection*{Текст задания}
Величина $Q=mv^2\sin^2\alpha$ должна быть не меньше $75$, если $m=3$, $v=10$, $0^\circ\le\alpha\le90^\circ$. Найдите минимальное значение удвоенного угла $2\alpha$ в градусах.

\subsection*{Ваше решение}
В работе записано:
\begin{align*}
75&=3\cdot100\cdot\sin^2\alpha,\\
75&=300\sin^2\alpha,\\
\sin^2\alpha&=\frac{75}{300},\\
\sin\alpha&=\sqrt{\frac{75}{300}},\\
\sin\alpha&=\sqrt{\frac5{10}}.
\end{align*}
Далее расчёт зачёркнут, итоговое значение $2\alpha$ не найдено.

\subsection*{Ваш ответ}
Не указан.

\subsection*{Ошибка}
\textcolor{red}{Ошибка:} неверное упрощение дроби под корнем, отсутствие объяснения перехода от неравенства к граничному равенству и незавершённый поиск $2\alpha$.

\subsection*{Где именно возникает ошибка}
Последняя надёжная правильная числовая запись:
\[
\sin\alpha=\sqrt{\frac{75}{300}}
\]
для граничного случая $Q=75$. Первый неверный числовой переход:
\[
\sqrt{\frac{75}{300}}\ne\sqrt{\frac5{10}}.
\]
Кроме того, исходное условие «не меньше $75$» сначала требует неравенства $Q\ge75$. Равенство $Q=75$ можно использовать при поиске минимального угла, но это необходимо обосновать.

\subsection*{Почему этот шаг неверен}
При сокращении дроби числитель и знаменатель делят на одно и то же ненулевое число. В данном случае
\[
\frac{75}{300}=\frac{75:75}{300:75}=\frac14,
\]
а $\frac5{10}=\frac12$. Эти дроби различны, поэтому и их квадратные корни различаются. Из-за неправильной дроби нельзя получить правильный порог для синуса.

Отдельно важно сохранить знак сравнения: формулировка «не меньше» означает $\ge$. На промежутке $0^\circ\le\alpha\le90^\circ$ функция $\sin\alpha$ неотрицательна и возрастает, а значит, вместе с ней возрастает $\sin^2\alpha$. Поэтому наименьший допустимый угол соответствует именно равенству $Q=75$.

\subsection*{Ключевая идея}
\textcolor{blue!60!black}{Ключевая идея:} сначала решить неравенство относительно $\sin^2\alpha$, затем учесть диапазон углов, найти наименьшее $\alpha$ и только после этого вычислить $2\alpha$.

\subsection*{Исправление}
В граничном случае вернитесь к правильной дроби $\frac{75}{300}$ и сократите её до $\frac14$. Тогда $\sin\alpha=\frac12$, минимальный угол равен $30^\circ$, а удвоенный --- $60^\circ$.

\subsection*{Полное правильное решение}
По условию
\[
Q=mv^2\sin^2\alpha\ge75.
\]
Подставим $m=3$ и $v=10$:
\begin{align*}
3\cdot10^2\sin^2\alpha&\ge75,\\
300\sin^2\alpha&\ge75,\\
\sin^2\alpha&\ge\frac{75}{300}=\frac14.
\end{align*}
Поскольку $0^\circ\le\alpha\le90^\circ$, синус неотрицателен. Следовательно, можно перейти от квадратов к неравенству
\[
\sin\alpha\ge\sqrt{\frac14}=\frac12.
\]
На промежутке от нуля до девяноста градусов синус возрастает. Так как $\sin30^\circ=\frac12$, получаем
\[
\alpha\ge30^\circ.
\]
Минимальный угол равен $30^\circ$. Требуется удвоенный угол, поэтому
\[
(2\alpha)_{\min}=2\cdot30^\circ=60^\circ.
\]

\subsection*{Проверка}
При $\alpha=30^\circ$ исходная величина равна
\[
Q=3\cdot10^2\cdot\sin^230^\circ
=300\cdot\frac14=75.
\]
Условие «не меньше $75$» выполнено точно на границе; при меньших углах в указанном диапазоне синус и $Q$ меньше.

\subsection*{Итог}
\textcolor{teal!60!black}{Итог:} $\boxed{(2\alpha)_{\min}=60^\circ}$.

\subsection*{Задача на закрепление}
Величина $Q=mv^2\sin^2\beta$ должна быть не меньше $100$, если $m=2$, $v=10$, $0^\circ\le\beta\le90^\circ$. Найдите минимальное значение угла $2\beta$ в градусах.

\clearpage
\phantomsection\label{task:8}
\section*{Задание № 8. Доля первой секунды}
\subsection*{Текст задания}
Скорость тела $v(t)=9\sin(\pi t)$ см/с. Какую долю первой секунды скорость была не меньше $4{,}5$ см/с? Округлите ответ до сотых.

\subsection*{Ваше решение}
Из записи видно, что неравенство и границы промежутка найдены правильно:
\begin{align*}
9\sin(\pi t)&\ge4{,}5,\\
\sin(\pi t)&\ge\frac12,\\
\frac\pi6&\le\pi t\le\frac{5\pi}6,\\
\frac16&\le t\le\frac56,\\
\frac56-\frac16&=\frac46=\frac23.
\end{align*}
После этого в ответе записано $0{,}62$.

\subsection*{Ваш ответ}
$0{,}62$.

\subsection*{Ошибка}
\textcolor{red}{Ошибка:} неверный перевод найденной правильной дроби $\frac23$ в десятичную запись и неверное округление до сотых.

\subsection*{Где именно возникает ошибка}
Последний правильный шаг:
\[
\frac56-\frac16=\frac23.
\]
Первый неправильный результат --- итог $0{,}62$. На самом деле $\frac23=0{,}666\ldots$, а не $0{,}62$.

\subsection*{Почему этот шаг неверен}
Дробь $\frac23$ означает деление двух на три. Получается бесконечная периодическая десятичная дробь: $0{,}666\ldots$. Чтобы округлить до сотых, оставляют две цифры после запятой и смотрят на третью: она равна шести. Поэтому цифру сотых увеличивают на единицу, получая $0{,}67$. Значение $0{,}62$ не является результатом такого округления.

\subsection*{Ключевая идея}
\textcolor{blue!60!black}{Ключевая идея:} если найден точный промежуток времени, сначала вычислить его длину дробью, затем разделить на продолжительность всей первой секунды и лишь в конце выполнить округление.

\subsection*{Исправление}
Сохраните правильный промежуток $\left[\frac16,\frac56\right]$ и результат $\frac23$; замените только последнюю десятичную запись:
\[
\frac23\approx0{,}67.
\]

\subsection*{Полное правильное решение}
Первая секунда соответствует $0\le t\le1$. Требование к скорости записывается как
\[
9\sin(\pi t)\ge4{,}5.
\]
Разделим обе части неравенства на положительное число девять:
\[
\sin(\pi t)\ge\frac{4{,}5}{9}=\frac12.
\]
За первую секунду аргумент $\pi t$ пробегает промежуток $[0,\pi]$. На нём синус равен $\frac12$ в точках $\frac\pi6$ и $\frac{5\pi}6$; между ними синус не меньше $\frac12$. Значит,
\[
\frac\pi6\le\pi t\le\frac{5\pi}6.
\]
Разделим всё неравенство на положительное число $\pi$:
\[
\frac16\le t\le\frac56.
\]
Продолжительность найденного промежутка:
\[
\Delta t=\frac56-\frac16=\frac46=\frac23\ \text{с}.
\]
Длина всей первой секунды --- одна секунда, поэтому искомая доля равна
\[
\frac{\Delta t}{1\ \text{с}}=\frac23=0{,}666\ldots\approx0{,}67.
\]

\subsection*{Проверка}
На концах найденного промежутка скорость равна $4{,}5$ см/с:
\[
9\sin\frac\pi6=9\sin\frac{5\pi}6=4{,}5.
\]
В середине, при $t=\frac12$, скорость равна $9$ см/с, то есть выше порога. Доля $\frac23$ находится между нулём и единицей, как и должна.

\subsection*{Итог}
\textcolor{teal!60!black}{Итог:} $\boxed{0{,}67}$ первой секунды.

\subsection*{Задача на закрепление}
Скорость тела $v(t)=10\sin(\pi t)$ см/с. Какую долю первой секунды скорость была не меньше $5\sqrt2$ см/с? Округлите ответ до сотых.

\vfill
\begin{center}
\small Лёвин Артём Александрович эксклюзивно для Ярослава Гаврилова
\end{center}
\end{document}
